\scanpage{105}
\mspage{8}
commute aux $\varprojlim$ de type $d$. Bien sûr, on a un
énoncé dual pour les $\varinjlim$

\uline{8.2.} Notons que pour \deleted{\unsure{T}} $\lambda$ quelconque \inserted{et un $\lambda$-type $T$}, il est vrai que pour
tt $\lambda$-foncteur \inserted{dans $\mathrm{Cat}_\lambda$}
\margin{(\unsure{remonte|remontez|remarque})}
\[
\deleted{\unsure{i:C\to C'}}\quad i:\deleted{\unsure{\xi}}C\longrightarrow C'\qquad\text{pl. fid. \inserted{(fid.)}}\tag*{(8.2.1)}
\]
le foncteur
\[
T(i):T(C)\longrightarrow T(C')\tag*{(8.2.2)}
\]
est pl fidèle \inserted{(fid.)}. Si $T$ est muni de $\unsure{b|\beta}:T\longrightarrow\tau^I$ \inserted{(``base'')}
\uline{fidèle}, i.e. $R=R_T^\lambda$ qui 2-représente $T$ est muni de
$\varphi:I\longrightarrow\mathrm{Ob}\,R_{\deleted{\unsure{T}}}$, alors \deleted{le carré} dans le cas pl. fidèle
\diagram{\begin{tikzcd}[ampersand replacement=\&]
 {(8.2.3)}\& T_{\lambda}(C)\arrow[r,hook]\arrow[d]\&T(C')\arrow[d]\\
 \& C^I\arrow[r,hook]\&C'{}^I
\end{tikzcd}}
est 2-cartésien [Cela signifie aussi, en termes de $R$,
que la s-catégorie pleine \deleted{de} $\underline{\mathrm{Hom}}_\lambda(R,C)$ de $\underline{\mathrm{Hom}}_\lambda(R,C')$
\deleted{s'identifie aux}\inserted{est formée des} $\lambda$-foncteurs $R\xrightarrow{\xi}C'$ tels que $\xi(\varphi(i))\in C$
pour tt $i\in I$ : en effet, comme $I$ \inserted{$\lambda$-}engendre $R$, il
s'ensuit que $\xi$ (qui commute aux lim de type $\lambda$)
envoie $R$ dans l'image essentielle de $C$.]
\inserted{En termes imagés, si $(E_i)_{i\in I}$ est une famille d'objets de $C$, ``il revient au $=$'' de se donner}
% slanted
\margin{sur ces derniers une structure de type $T$ \deleted{\unsure{dans}} sur \unsure{$C'$|$E_i$} \deleted{\unsure{fon}} et ceci pour les morphismes}

Appliquant ceci, pour $C\in\mathrm{Ob}\,\mathrm{Cat}_\lambda$, au plongt
canonique
\[
i:C\lhook\joinrel\longrightarrow\widehat C\qquad\text{qui est un $\lambda$-foncteur si }\uline{$\lambda=(\underset{\leftarrow}{\lambda},\varnothing)$}\ ;\tag*{(8.2.4)}
\]
on trouve \deleted{\unsure{que}} le foncteur
\diagram{\begin{tikzcd}[ampersand replacement=\&]
 {(8.2.5)}\& T(C)\arrow[r,dash]\arrow[d]\&T(\widehat C)\arrow[d]\\
 \&C^I\arrow[r,dash]\&\widehat C{}^I
\end{tikzcd}}
(2-cartésien)
\margin{\deleted{sup.} On suppose $\underset{\rightarrow}{\lambda}=\varnothing$}

% slanted
\margin{\uline{NB} L'hypothèse $\widehat C$ une $\mathcal U$-catégorie n'est pas essentielle en fait, car (8.2.5) ou la représentabilité de $T_\lambda$ donne à $T_\lambda(C)$ pour tte $\lambda$-catégorie $C$ (pas néc. $\in\mathrm{Ob}\,\mathrm{Cat}_\lambda$) un sens à l'équiv. près, (....) si \deleted{\unsure{les}} $T(\widehat C)$ est connu. On va essayer de le ``calculer''.}
% slanted
\margin{On suppose désormais $\lambda$ quelc.}

\deleted{On a donc pour \unsure{cela|ici},} \inserted{On suppose ici} $\widehat C$ \deleted{est une $\lambda$-catégorie} \inserted{$\in\mathrm{Ob}\,\mathrm{Cat}_{\lambda}$} \deleted{(ceci suppose que c'est} \inserted{\deleted{\unsure{soit}} que $\mathrm{ss}\,\widehat C$ est} une $\mathcal U$-catégorie \inserted{(vrai p.ex. si $C$ est $\mathcal U$-petite)} \deleted{\unsure{est de même} \unsure{comme} une \unsure{cat|catégorie}}. Donc $T(C)$ est
\deleted{On a pour tt $X\in\mathrm{Ob}$} Considérons le foncteur
contravariant canonique de $\widehat C$ dans $\mathrm{Hom}(C,(\mathrm{Ens}))$, se
restreignant à $C$ \deleted{\unsure{au sens}} moyennant $i$ est un foncteur
canonique $C^o\longrightarrow\underline{\mathrm{Hom}}_\lambda(\widehat C,(\mathrm{Ens}))$, et qui s'en va en
fait dans la \deleted{$\unsure{C\longrightarrow|\widehat C}$} ss-catégorie pleine des foncteurs
qui commutent à tt type \inserted{``petit''} de lim (ind ou proj),
donc en particulier
\[
 C^o\longrightarrow\underline{\mathrm{Hom}}_\lambda(\widehat C,(\mathrm{Ens}))\qquad X\longmapsto(F\longmapsto F(X))\tag*{(8.2.6)}
\]
Or on a
\[
\underline{\mathrm{Hom}}_\lambda(\widehat C,(\mathrm{Ens}))\longrightarrow
\underline{\mathrm{Hom}}_{\unsure{\lambda|+}}(T_{\deleted{\unsure{\lambda}}}(\widehat C),T_{\deleted{\unsure{\lambda}}}(\mathrm{Ens}))
\]
% archive sheet number: 104
